% TOP_PROBLEM: {"id": 8900008, "title": "Congruent number decision problem", "category_id": 3, "category": {"id": 3, "name": "graph_theory", "display_name": "Graph Theory", "description": "Problems involving graphs, networks, and their properties.", "slug": "graph-theory", "order_index": 3, "created_at": "2026-07-31T15:26:25.671Z"}, "status": "open"}
% FIRST_POSTED: null
% !TeX program = lualatex
% Compile twice: lualatex 8900008.tex
% All references and prose are embedded; no BibTeX database or external inputs.
\documentclass[11pt,a4paper]{article}
\usepackage[margin=24mm,headheight=14pt]{geometry}
\usepackage{fontspec}
\setmainfont{Latin Modern Roman}
\setsansfont{Latin Modern Sans}
\setmonofont{Latin Modern Mono}
\usepackage{amsmath,amssymb,mathtools}
\usepackage{unicode-math}
\setmathfont{Latin Modern Math}
\usepackage{microtype}
\usepackage{enumitem,xurl,needspace}
\usepackage[dvipsnames]{xcolor}
\usepackage{hyperref}
\hypersetup{unicode=true,colorlinks=true,linkcolor=MidnightBlue,urlcolor=MidnightBlue,
 pdftitle={Research attempt: Problem 8900008},pdfauthor={Research notebook},bookmarksnumbered=true}
\usepackage{bookmark}
\usepackage{tocloft}
\setlength{\cftsecnumwidth}{3em}
\cftsetpnumwidth{2.5em}
\cftsetrmarg{4em}
\usepackage{titlesec}
\titleformat{\section}{\Large\bfseries\raggedright}{\thesection}{0.7em}{}
\usepackage{fancyhdr}
\pagestyle{fancy}\fancyhf{}
\fancyhead[L]{\small\sffamily Open Problems Math | Research notebook}
\fancyhead[R]{\small Problem 8900008 | 27 September 2026}
\fancyfoot[C]{\thepage}
\setlength{\parindent}{0pt}
\setlength{\parskip}{3pt plus 1pt}
\setlength{\emergencystretch}{3em}
\setcounter{tocdepth}{1}
\setcounter{secnumdepth}{1}
\setlist[itemize]{leftmargin=3em,itemsep=5pt,topsep=4pt,parsep=0pt}
\allowdisplaybreaks
\newcommand{\N}{\mathbb{N}}
\newcommand{\Z}{\mathbb{Z}}
\newcommand{\Q}{\mathbb{Q}}
\newcommand{\R}{\mathbb{R}}
\newcommand{\C}{\mathbb{C}}
\newcommand{\F}{\mathbb{F}}
\newcommand{\A}{\mathbb{A}}
\newcommand{\PP}{\mathbb P}
\newcommand{\E}{\mathbb{E}}
\newcommand{\eps}{\varepsilon}
\newcommand{\dd}{\,\mathrm d}
\newcommand{\sref}[2]{\hyperlink{source:#1:#2}{[S#2]}}
\newcommand{\eref}[2]{\hyperlink{extra:#1:#2}{[E#2]}}
% Turn the next switch off for a shorter reading copy. The quotations preserve
% every exactTarget field, including qualifications omitted from a short summary.
\newif\ifshowcatalogue
\showcataloguetrue
\newcommand{\cataloguescope}[1]{\ifshowcatalogue
 \paragraph{Exact scope in the supplied catalogue.}
 \begin{quote}\small #1\end{quote}\fi}
\usepackage{etoolbox}
\AtBeginEnvironment{itemize}{\small}
\numberwithin{equation}{section}
% Standalone export. Original research content preserved.
\begin{document}
\setcounter{section}{0}
% ================================================================
% problemId: 8900008
% Canonical problem ID: 8900008
% exactTarget SHA-256: c33ccbbc30631bf5dd7d361121550db41aec74611cfd1bc5b1c56274a5a496cf
\clearpage
\section*{\texorpdfstring{Congruent number decision problem}{Congruent number decision problem}}\label{8900008}
\subsection{Definitions and mathematical statement}
A positive integer $n$ is congruent if there exist $a,b,c\in\Q_{>0}$ with
\[
 a^2+b^2=c^2,\qquad ab/2=n.
\]
Equivalently, the elliptic curve $E_n:y^2=x^3-n^2x$ has a rational point with $y\ne0$. For such a point one obtains side lengths $|x^2-n^2|/|y|$, $2n|x|/|y|$, and $(x^2+n^2)/|y|$. Multiplication of $n$ by a rational square does not change congruence, so positive integers reduce to squarefree parts. The problem seeks an unconditional exact characterization, not only a conditional criterion.
\par\smallskip\textit{Formulation and concept references:} \sref{8900008}{1}.
\cataloguescope{Determine exactly which positive integers are congruent numbers, i.e. areas of right triangles with all three sides rational, equivalently for which squarefree n the elliptic curve y\textasciicircum{}2 = x\textasciicircum{}3 - n\textasciicircum{}2 x has a rational point with y != 0.}
\subsection{Short English statement}
Which positive integers occur as the areas of right triangles with entirely rational side lengths?
% BEGIN RESEARCH ADDITION 146
\subsection{Research attempt}

\paragraph{Current frontier.}
Tunnell supplies an unconditional obstruction and, conditional on the weak
Birch--Swinnerton-Dyer conjecture for the relevant curves, a necessary and
sufficient test. For odd squarefree $n$, define
\[
 f(n)=\#\{(x,y,z)\in\Z^3:x^2+2y^2+8z^2=n\},\qquad
 g(n)=\#\{(x,y,z)\in\Z^3:x^2+2y^2+32z^2=n\}.
\]
Congruence implies $f(n)=2g(n)$; the reverse implication is the conditional
part. The even test is also in Conrad's Theorem 4.11. \sref{8900008}{1}\eref{8900008}{1}
The 2026 published rank-zero $p$-converse for CM elliptic curves, combined
with Selmer distributions, proves a density-$1/2$ analytic-rank-zero assertion
in the positive squarefree congruent-number twist family, not an exact
every-integer decision. \eref{8900008}{3} Tian separates the universal,
statistical and Heegner-point questions explicitly. \eref{8900008}{2}

\paragraph{Attack route.}
A square-class descent allocates the primes dividing $n$ among four factors.
The tempting missing lemma is that local solubility of one of the resulting
systems ensures global solubility. We derive the finite allocation exactly
and test this lemma on an explicit member of that same family. Proving the
converse to Tunnell's equality would be another route, but replacing a
vanishing $L$-value by positive rank without an additional theorem would
assume the central difficulty.

\paragraph{Attempt: a finite allocation, not yet a finite decision.}
For each positive squarefree $n$, consider the $4^{\omega(n)}$ ordered
quadruples of pairwise coprime positive squarefree integers $a,b,c,d$ with
$abcd=n$. We claim that $n$ is congruent exactly when at least one system
\begin{equation}\label{8900008:descent}
 ar^2-bs^2=ct^2,\qquad ar^2+bs^2=du^2
\end{equation}
has positive rational $r,s,t,u$.

A rational right triangle is a positive rational multiple of a primitive
integer triple $(m^2-v^2,2mv,m^2+v^2)$, where $m>v$, $\gcd(m,v)=1$ and
$m,v$ have opposite parities. The parametrization and square-class viewpoint
are recalled in Tian, Section 1. \eref{8900008}{2} The four integers
$m,v,m-v,m+v$ are pairwise coprime: a common divisor of the last two divides
$2$, and both are odd. Their squarefree parts $a,b,c,d$ multiply to the
squarefree part $n$ of the area $mv(m-v)(m+v)$. Their square parts give
\eqref{8900008:descent}, proving necessity.

Conversely, a positive rational solution gives $m=ar^2$, $v=bs^2$, $m>v>0$
and
\[
 mv(m^2-v^2)=abcd(rstu)^2=n(rstu)^2.
\]
Therefore
\begin{equation}\label{8900008:triangle}
 \frac{m^2-v^2}{rstu},\quad \frac{2mv}{rstu},\quad
 \frac{m^2+v^2}{rstu}
\end{equation}
are positive rational sides of a right triangle of area $n$. This proves
the equivalence, but each remaining rational-point question can still be
hard; no height bound or terminating search has appeared.

\paragraph{An everywhere locally soluble obstruction inside this descent.}
For $(a,b,c,d)=(17,1,1,1)$, consider
\begin{equation}\label{8900008:C}
 C:\quad 17r^2-s^2=t^2,\qquad17r^2+s^2=u^2
 \quad\subset\PP^3_{\Q}.
\end{equation}
We prove that $C$ has points over every completion of $\Q$ but none over
$\Q$. This disproves the proposed local-to-global lemma, not a proposed
classification of all congruent numbers.

Over characteristic different from $2,17$, the Jacobian rows are
\[
 (34r,-2s,-2t,0),\qquad(34r,2s,0,-2u).
\]
A nonzero linear dependence with coefficients $\lambda,\mu$ implies
$(\lambda+\mu)r=(-\lambda+\mu)s=\lambda t=\mu u=0$.
If one coefficient vanishes, the defining equations force every coordinate
zero. If both are nonzero, $t=u=0$, and adding and subtracting the equations
again forces $r=s=0$. Thus the intersection is smooth. A smooth complete
intersection of two quadrics in $\PP^3$ is geometrically connected, and
adjunction gives canonical degree zero, hence genus one.
For $p\ne2,17$, Hasse--Weil therefore gives
$\#C(\F_p)\ge p+1-2\sqrt p>0$, and smoothness lifts a point to $\Q_p$.
This uses the genus-one case of Poonen's Corollary 7.2.1. \eref{8900008}{4}

At $p=2$, $17\equiv1\pmod8$ is a square in $\Q_2$, yielding
$[1:0:\sqrt{17}:\sqrt{17}]$. At $p=17$, $-1$ is a square and
$[0:1:\sqrt{-1}:1]$ works. Over $\R$, use $[1:1:4:3\sqrt2]$.
There are even local points with no zero coordinates: a smooth local curve
has infinitely many points near a local point, while its intersection with
the coordinate hyperplanes is finite. The displayed real point is already
positive.

No rational point on $C$ can have a zero coordinate. If $r=0$ or $u=0$,
a sum of rational squares forces every coordinate zero. If $s=0$ or $t=0$,
a remaining equation makes $17$ a rational square unless all coordinates
vanish. Thus any rational point could be made positive by changing signs,
and would imply congruence of $17$ via \eqref{8900008:triangle}. However exact
enumeration gives
\[
 f(17)=16,\qquad g(17)=4.
\]
Tunnell's unconditional necessary condition excludes $17$, proving
$C(\Q)=\varnothing$. \sref{8900008}{1}\eref{8900008}{1} The four-factor allocation
has therefore not removed the Hasse-principle obstruction.

\paragraph{Stress test and source audit.}
This countertest lies inside the derived family, not an unrelated genus-one
example. But it does not refute a route first imposing Tunnell's equality:
$17$ fails that filter. The remaining assertion there would be that every
integer passing the appropriate equality has a rational point on at least
one descent curve. No proof is provided.

A 2024 manuscript found during current research is not adopted as a solution.
Its Theorem 1.1 attributes to Smith both rank one and $L(E_n,1)\ne0$ for
all $n\equiv5,6,7\pmod8$. \eref{8900008}{5} In the squarefree cases the root
number is $-1$, forcing $L(E_n,1)=0$ by the functional equation, as Tian
records in Section 1. \eref{8900008}{2} That displayed assertion therefore
cannot supply a universal theorem. A statistical Selmer result must not be
turned into an every-integer claim.

The script checks Tunnell's finite counts, the allocation and triangle
identities on primitive triples, and smooth points modulo small good primes.
The assertion for every prime rests on Hasse--Weil, and rational nonexistence
on Tunnell's obstruction, not on any bounded search.

\paragraph{Outcome.}
\textbf{Outcome: Exploratory approach with unresolved gap.}
The exact square-class allocation and an explicit failure of the local
solubility shortcut have been proved. This combines elementary descent with
established theorems and makes no novelty claim. It gives no unconditional
complete classification, no uniform height bound for a first triangle and
no proof of the converse to Tunnell's test.

% END RESEARCH ADDITION 146
\subsection{Sources}
\begin{itemize}\raggedright
\item[\textnormal{[S1]}]\hypertarget{source:8900008:1}{} Keith Conrad, "The Congruent Number Problem," expository notes, University of Connecticut, kconrad.math.uconn.edu/blurbs/ugradnumthy/congnumber.pdf.
\url{https://kconrad.math.uconn.edu/blurbs/ugradnumthy/congnumber.pdf}.
{\small\textit{Catalogue formulation source.}}
% BEGIN RESEARCH REFERENCES 146
\item[\textnormal{[E1]}]\hypertarget{extra:8900008:1}{} J. B. Tunnell,
\emph{A classical Diophantine problem and modular forms of weight $3/2$},
Inventiones Mathematicae 72 (1983), 323--334, introductory theorem and Section 3.
\url{https://sites.math.rutgers.edu/~zeilberg/EM22/JT1983.pdf}.
\item[\textnormal{[E2]}]\hypertarget{extra:8900008:2}{} Y. Tian,
\emph{The congruent number problem and elliptic curves}, ICM 2022 Proceedings,
Vol. 3, 1990--2010, Section 1.
\url{https://ems.press/content/book-chapter-files/33187}.
\item[\textnormal{[E3]}]\hypertarget{extra:8900008:3}{} A. A. Burungale and Y. Tian,
\emph{A rank zero $p$-converse to a theorem of Gross--Zagier, Kolyvagin and
Rubin}, Annals of Mathematics 203 (2026), 1--13, main theorem and Goldfeld
application. \url{https://annals.math.princeton.edu/2026/203-1/p01}.
\item[\textnormal{[E4]}]\hypertarget{extra:8900008:4}{} B. Poonen,
\emph{Rational Points on Varieties}, GSM 186, AMS, 2017, Corollary 7.2.1.
\url{https://math.mit.edu/~poonen/papers/Qpoints.pdf}.
\item[\textnormal{[E5]}]\hypertarget{extra:8900008:5}{} H. Chen, R. Ma and T. Du,
\emph{A exact way to verify whether n is a congruent number using Heegner
points}, arXiv:2411.14188v1, Theorem 1.1.
\url{https://arxiv.org/html/2411.14188v1}.
{\small\textit{Consulted to audit its displayed claim, not used as a valid
theorem; see the contradiction in the stress test.}}

% END RESEARCH REFERENCES 146
\end{itemize}

\end{document}
